% lagrangian-product-algorithm.tex — \input'd from algorithms.tex
% Conventions: thesis/CLAUDE.md
%
% Sources: Bezdek-Bezdek (2009), Rudolf (2022), Artstein-Avidan & Ostrover (2014),
%   HK2017. Agent-written, math verified by Joern.
%
% Structure: The Billiard Algorithm

% Jörn: math approved (e1345bb) — from \subsection{The Billiard Algorithm} to end of subsection

\begin{theorem}[Completeness of the billiard enumeration]%
\label{thm:billiard-completeness}
Let $K = K_q \times K_p$ be a Lagrangian product.
Then $c_{\mathrm{EHZ}}(K)$ equals the minimum of $A(S, \sigma)$
(equation~\eqref{eq:action-formula})
over all $(S, \sigma)$ with $\sigma$ matching the pattern
$\bigl([Q \mid QQ]\,[P \mid PP]\bigr)^k$
for $k \in \{2, 3\}$.
\end{theorem}

\begin{proof}
\textbf{Overview.}\quad
Every valid $(S, \sigma)$ from the restricted set is also a valid
input to the algorithm of Section~\ref{sec:main-result}, so
\begin{equation}\label{eq:billiard-lb}
  \min_{\text{restricted}} A(S, \sigma)
  \geq \min_{\text{all}} A(S, \sigma)
  = c_{\mathrm{EHZ}}(K).
\end{equation}
It remains to show the reverse inequality.

\medskip\noindent
\textbf{Step 1: simple minimizer has alternating block structure.}\quad
By Theorem~\ref{thm:simple-minimizer}, $c_{\mathrm{EHZ}}(K)$ is
achieved by a simple Reeb orbit $\gamma^*$ with representation
$(\sigma^*, \tau^*)$, all $\tau_i > 0$.
By Lemma~\ref{lem:sigma-structure}, $\sigma^*$ has the form
$\bigl([Q \mid QQ]\,[P \mid PP]\bigr)^{k^*}$ for some $k^* \geq 2$
(Lemma~\ref{lem:k1-infeasible} excludes $k^* = 1$).
Hence
\begin{equation}\label{eq:billiard-ub1}
  \min_{\text{alternating-block}} A(S, \sigma) \leq A(S^*, \sigma^*)
  = c_{\mathrm{EHZ}}(K).
\end{equation}

\medskip\noindent
\textbf{Step 2: bound $k^* \leq 3$.}\quad
By the billiard characterization
(Theorem~\ref{thm:billiard-characterization}), $c_{\mathrm{EHZ}}(K)$
is also the minimum $K_p^\circ$-length over billiard trajectories
in~$K_q$. By the bounce bound (Theorem~\ref{thm:bounce-bound}),
this minimum is achieved by a trajectory with at most $3$ bounce points.
Each bounce point corresponds to one $q$-block in the simple orbit
(Lemma~\ref{lem:sigma-structure}), so $k^* \leq 3$.

\medskip\noindent
\textbf{Conclusion.}\quad
Combining~\eqref{eq:billiard-lb} and~\eqref{eq:billiard-ub1}
with the restriction $k \in \{2, 3\}$:
\[
  c_{\mathrm{EHZ}}(K)
  \leq \min_{k \in \{2,3\}} A(S, \sigma)
  \leq \min_{\text{alternating-block}} A(S, \sigma)
  \leq c_{\mathrm{EHZ}}(K). \qedhere
\]
\end{proof}

\begin{remark}[Complexity]\label{rem:billiard-complexity}
% Downstream: total iterations for k=2,3 combined.
%   For each k: choose k q-blocks from O(n_q) blocks, k p-blocks
%   from O(n_p) blocks, then enumerate orderings.
%   Dominated by k=3: O(n_q^3 * n_p^3) iterations.
%   Each iteration: solve (|sigma|+5) x (|sigma|+5) system with
%   |sigma| <= 12, so each solve is O(1).
Let $n_q = F_q$ and $n_p = F_p$ denote the edge counts of $K_q$
and $K_p$.
The number of $(S, \sigma)$ to evaluate is
$O(n_q^3 \, n_p^3)$, dominated by $k = 3$:
\begin{itemize}
\item Choose $3$ non-overlapping $q$-blocks: $O(n_q^3)$ choices.
\item Choose $3$ non-overlapping $p$-blocks: $O(n_p^3)$ choices.
\item Orderings per selection: at most $2! \times 3! = 12$
  (fix first $q$-block to remove cyclic symmetry,
  permute remaining $q$-blocks and all $p$-blocks).
\item Within-block orderings: at most $2^6 = 64$
  (each of 6 blocks can have reversed internal order).
\end{itemize}
Each evaluation solves a linear system of size at most
$12 + 5 = 17$, which is $O(1)$.
The total cost is polynomial in $n_q$ and $n_p$.
\end{remark}

\medskip
\noindent\textbf{Algorithm}\label{alg:billiard}
(computes $c_{\mathrm{EHZ}}(K_q \times K_p)$
via Theorem~\ref{thm:billiard-completeness}).

\smallskip\noindent
\emph{Input:} Lagrangian product $K = K_q \times K_p$ in
$H$-representation.
\quad
\emph{Output:} $c_{\mathrm{EHZ}}(K)$.

\begin{enumerate}
\item \textbf{Classify} facets into $q$-type and $p$-type
  (Lemma~\ref{lem:lagrangian-facets}).
\item \textbf{Build blocks.}
  For each polygon ($K_q$ and $K_p$), enumerate all
  \emph{single-facet blocks} (one edge) and
  \emph{pair blocks} (two adjacent edges, both orderings).
\item \textbf{Enumerate.}
  For $k = 2, 3$: for each selection of $k$ non-overlapping
  $q$-blocks and $k$ non-overlapping $p$-blocks, and each
  interleaving into the pattern
  $\bigl([Q \mid QQ]\,[P \mid PP]\bigr)^k$,
  flatten the blocks into a facet sequence~$\sigma$.
\item \textbf{Solve and filter.}
  For each $\sigma$: solve the KKT system~\eqref{eq:linear-system},
  filter by $\beta_i > 0$ and $Q(\beta) > 0$, compute
  $A(S, \sigma) = \tfrac{1}{2}\, Q(\beta)^{-1}$.
\item \textbf{Return} $\min A(S, \sigma)$.
\end{enumerate}
